\section{Motivación y pertinencia}
\label{sec:motivacion_pertinencia}

Para mejorar esta situación, en 2025 se probó una primera página web de inscripción. Sin embargo, esa página tenía un problema muy grande: no contaba con una opción para que el estudiante recuperara su contraseña en caso de olvido. Por eso, muchos jóvenes olvidaron sus datos de ingreso y tuvieron que volver al Módulo para solicitar la restitución manual de su cuenta. Mientras realizaban ese trámite, los cupos de las carreras más buscadas se terminaban y varios tuvieron que inscribirse en otra especialidad que no era la que deseaban.

A esto se sumó otro problema: la página solo permitía realizar la preinscripción al cierre de una gestión, pero no permitía completar la inscripción definitiva al inicio de la siguiente. En el BTH, para completar la inscripción se deben presentar las libretas de 1ro a 3ro de secundaria, la cédula de los tutores, la fotocopia del carnet del estudiante y su certificado de nacimiento. Como no había un espacio en el sistema para registrar y revisar esos documentos, los expedientes quedaban incompletos, la documentación se desordenaba y el personal enfrentaba una sobrecarga de trabajo.

Tampoco existía una forma automática de distribuir a los estudiantes en cursos y paralelos. En el BTH, primero el estudiante elige una especialidad y luego la dirección debe ubicarlo en un curso y paralelo según el colegio del que proviene. Al realizar este reparto de forma manual, en planillas aparte, aparecían errores, algunos cursos quedaban muy llenos y otros vacíos, y se producían cruces de horarios entre el Módulo y los colegios.

La motivación principal nace de lo vivido por la propia comunidad de San Julián. Observar que, por olvidar una contraseña o por hacer largas filas, un estudiante podía perder la oportunidad de estudiar la carrera que deseaba, y ver al personal revisando carpetas a mano y elaborando listas con mucho esfuerzo, impulsó al equipo de la especialidad de Sistemas Informáticos a buscar una solución completa y sencilla.

Este trabajo también cuenta con respaldo legal y académico. El Reglamento de Graduación del Bachillerato Técnico Humanístico, aprobado por la Resolución Ministerial RM 0912/2023 \parencite{ministerioeducacion2023}, establece que los proyectos deben resolver problemas reales de la institución aplicando lo aprendido en la especialidad. En la misma línea, \textcite{unesco2024} plantea que la transformación digital educativa requiere coordinación, capacidad y cultura institucional, y \textcite{unesco2023} subraya que la tecnología debe gestionarse con el estudiante en el centro, lo que justifica desarrollar una página que no solo reciba datos, sino que organice todo el proceso de inscripción.
